\clearpage\phantomsection\label{sujet:9-1}\addcontentsline{toc}{chapter}{Final mock · Practice 1}

\section*{Final mock · Practice 1}

{\sffamily\bfseries\color{teal}PRACTICE PAPER · TRAINING\par}

\textbf{Duration: 4 h · Marked out of 20 · Four independent exercises.}\\Calculator allowed. Every answer must be justified; incomplete working will be taken into account. Give exact values; round probabilities to $10^{-6}$ and other approximations to $10^{-4}$, unless stated otherwise.

\smallskip\textbf{Name:}\enspace\dotfill\quad\textbf{Class:}\enspace\dotfill\quad\textbf{Date:}\enspace\dotfill

\subsection*{Exercise 1 · A logarithmic function \hfill 5 points}Consider the function $f$ defined on $]0;+\infty[$ by $f(x)=\frac{2+\ln x}x$. Let $\mathcal C_f$ denote its graph.

\begin{enumerate}

\item Determine the limits of $f$ at $0^+$ and at $+\infty$. \hfill\textit{(0.5 pt)}

\item Calculate $f' (x)$, then draw up the table of variations of $f$. \hfill\textit{(0.5 pt)}

\item State the maximum value of $f$, then solve $f(x)=0$. \hfill\textit{(0.5 pt)}

\item Study the convexity of $f$ and determine the coordinates of the point of inflection of $\mathcal C_f$. \hfill\textit{(1 pt)}

\item Prove that the equation $f(x)=\frac12 e^{1}$ has exactly two solutions in $]0;+\infty[$. \hfill\textit{(0.5 pt)}

\item Find an antiderivative of $f$ on $]0;+\infty[$. \hfill\textit{(0.5 pt)}

\item Calculate the mean value of $f$ on $[1;e]$. \hfill\textit{(0.5 pt)}

\item Independently of the previous questions, solve $\ln(x-2)+\ln(x+2)=\ln(32)$, after stating the set on which this equation is defined. \hfill\textit{(1 pt)}

\end{enumerate}

\subsection*{Exercise 2 · Periodic phenomenon \hfill 5 points}A point moves on a wheel turning at constant speed in a vertical plane. Its height above the ground is modelled by $h(t)=9-3\cos\left(\frac{\pi t}{5}\right)$, where $t$ is expressed in seconds and $h(t)$ in metres. We study one complete revolution, for $t\in[0;10]$.

\begin{enumerate}

\item Give a period of the model extended to $\R$, then the minimum and maximum heights over one revolution. \hfill\textit{(0.5 pt)}

\item Calculate $h'(t)$ and draw up the table of variations of $h$ on $[0;10]$. \hfill\textit{(0.5 pt)}

\item Find all times $t\in[0;10]$ such that $h(t)=\frac{21}{2}$. \hfill\textit{(0.5 pt)}

\item Calculate the mean value of $h$ on $[0;10]$. \hfill\textit{(1 pt)}

\item Study the convexity of $h$ on the interval of study and determine the coordinates of the points of inflection of its graph. \hfill\textit{(0.5 pt)}

\item Determine the total time during which the height is at least equal to the level in question 3. \hfill\textit{(0.5 pt)}

\item The vertical speed is given by $h'(t)$. Determine its maximum and minimum values, in metres per second, and the times at which they are reached. \hfill\textit{(0.5 pt)}

\item Calculate the vertical displacement $h(10)-h(0)$, then the sum of the height travelled on the way up and that travelled on the way down. Explain why these two results are different. \hfill\textit{(1 pt)}

\end{enumerate}

\clearpage\noindent{\sffamily\bfseries\color{navy}Final mock · Practice 1 · continued}\par

\subsection*{Exercise 3 · Line, plane and projection \hfill 5 points}In an orthonormal coordinate system of space, consider the plane $\mathcal P$ with equation $x+2y+2z=11$ and the point $M(2;0;0)$. The line $d$ passes through $M$ and has direction vector $\vec n=(1;2;2)$.

\begin{enumerate}

\item Give a parametric representation of $d$. \hfill\textit{(0.5 pt)}

\item Prove that $d$ meets $\mathcal P$ at a point $H$, then calculate the coordinates of $H$. \hfill\textit{(0.5 pt)}

\item Justify that $H$ is the orthogonal projection of $M$ onto $\mathcal P$, then calculate the distance from $M$ to $\mathcal P$. \hfill\textit{(0.5 pt)}

\item Consider the plane $\mathcal Q$ with equation $y=z$. Prove that $\mathcal P\cap\mathcal Q$ is a line and give a parametric representation of it. \hfill\textit{(1 pt)}

\item Check that $H\in\mathcal Q$. Deduce the distance from $M$ to the line $\mathcal P\cap\mathcal Q$. \hfill\textit{(0.5 pt)}

\item Determine a Cartesian equation of the plane parallel to $\mathcal P$ passing through $M$. \hfill\textit{(0.5 pt)}

\item Determine the angle between $d$ and the line $\mathcal P\cap\mathcal Q$. \hfill\textit{(0.5 pt)}

\item Define the point $N$ by $\vect{MN}=\vec n$. Calculate the distances from $N$ to the plane $\mathcal P$ and to the line $\mathcal P\cap\mathcal Q$, then compare them. \hfill\textit{(1 pt)}

\end{enumerate}

\subsection*{Exercise 4 · Repetitions and estimation \hfill 5 points}We repeat $20$ times a trial with two outcomes, called success and failure. The trials are independent and the probability of success is $p=\frac{1}{5}$ on each trial. Let $X$ denote the number of successes obtained.

\begin{enumerate}

\item Justify that $X$ follows a binomial distribution and state its parameters. \hfill\textit{(0.5 pt)}

\item Calculate the probability of obtaining exactly two successes. \hfill\textit{(0.5 pt)}

\item Using complementary events, express and then calculate the probability of obtaining at least two successes. \hfill\textit{(0.5 pt)}

\item Calculate the expected value and the variance of $X$. \hfill\textit{(1 pt)}

\item We now repeat the same trial a number $N\ge1$ of times, still independently. Determine the smallest integer $N$ for which the probability of obtaining at least one success is greater than or equal to $0.95$. A search using a table or a loop is allowed. \hfill\textit{(0.5 pt)}

\item For an integer $j\in\{0,\ldots,20\}$, prove the formula for $P(X=j)$ by counting the possible positions of the successes. \hfill\textit{(0.5 pt)}

\item For $N$ independent trials with the same parameter $p$, let $F_N$ denote the frequency of successes. Prove the expressions for $E(F_N)$ and $V(F_N)$ in terms of $N$ and $p$. \hfill\textit{(0.5 pt)}

\item Using the Bienaymé–Chebyshev inequality, determine a sufficient size $N$ so that $P(|F_N-p|<0.02)\ge0.95$. Is the size obtained necessarily the smallest possible for this probability? Justify. \hfill\textit{(1 pt)}

\end{enumerate}

\ifprof\clearpage

\section*{Solutions · Final mock · Practice 1}\label{cor:9-1}

\subsection*{Exercise 1 · A logarithmic function}\begin{enumerate}

\item \textit{[0.5 pt]} At $0^+$, the numerator tends to $-\infty$ and the denominator to $0^+$: limit $-\infty$; at $+\infty$, limit $0$ by comparison of growth rates.

\item \textit{[0.5 pt]} By the quotient rule,
\[f'(x)=\frac{(\frac{1}{x})x-(2+\ln x)}{x^2}=\frac{-1-\ln x}{x^2}.\]
The denominator is positive. The numerator vanishes at $s=e^{-1}$, is positive for $0<x<s$ and negative for $x>s$. Hence $f$ is strictly increasing on $]0,s]$, then strictly decreasing on $[s,+\infty[$.

\item \textit{[0.5 pt]} The maximum is attained only at $s=e^{-1}$ and equals
\[f(s)=\frac{2- 1}{e^{-1}}=e^{1}.\]
Since $x>0$, $f(x)=0$ is equivalent to $2+\ln x=0$, that is $x=e^{-2}$.

\item \textit{[1 pt]} Differentiating $(-1-\ln x)x^{-2}$:
\[f''(x)=-x^{-3}-2(-1-\ln x)x^{-3}=\frac{2\ln x+1}{x^3}.\]
The numerator vanishes at $t=e^{-\frac{1}{2}}$ and changes from negative to positive. The function is concave on $]0,t]$, convex on $[t,+\infty[$. It changes convexity at $t$ and $f(t)=\frac{3}{2t}$: the point of inflection is $I(t,\frac{3}{2t})$.

\item \textit{[0.5 pt]} Let $m=\frac{e^{1}}{2}$ and $s=e^{-1}$. The function is continuous on $]0,+\infty[$. On $]0,s]$, it is strictly increasing, with limit $-\infty$ at $0$ and value $2m$ at $s$: the corollary of the IVT gives exactly one solution in $]0,s[$. On $[s,+\infty[$, it is strictly decreasing from $2m$ to $0<m$: exactly one solution in $]s,+\infty[$. These intervals cover the domain and $f(s)\ne m$: there are exactly two solutions.

\item \textit{[0.5 pt]} Let $F(x)=2\ln x+\frac12(\ln x)^2$, differentiable for $x>0$. Then $F'(x)=\frac{2}{x}+\frac{\ln x}{x}=f(x)$; $F$ is an antiderivative on $]0,+\infty[$.

\item \textit{[0.5 pt]} With this antiderivative, $F(e)=\frac{5}{2}$ and $F(1)=0$. The integral therefore equals $\frac{5}{2}$, and the mean value is
\[\mu=\frac{1}{e-1}\int_1^e f(x)\,\dd x=\frac{5}{2(e-1)}.\]

\item \textit{[1 pt]} Domain $x>2$; $x^2-4=32$, so $x=6$, the negative root being excluded.

\end{enumerate}

\subsection*{Exercise 2 · Periodic phenomenon}\begin{enumerate}

\item \textit{[0.5 pt]} Period $10$ s; minimum $6$ m, maximum $12$ m.

\item \textit{[0.5 pt]} $h'(t)=(\frac{3}{5}\pi)\sin (\frac{\pi t}{5})$; increasing on $[0,5]$, decreasing on $[5,10]$.

\item \textit{[0.5 pt]} The equation is equivalent to $\cos (\frac{\pi t}{5})=-\frac{1}{2}$. Over one revolution, $\frac{\pi t}{5}\in[0,2\pi]$; the suitable angles are $\frac{2\pi}{3}$ and $\frac{4\pi}{3}$. Hence $t=\frac{10}{3}$ s or $t=\frac{20}{3}$ s.

\item \textit{[1 pt]} An antiderivative is $F(t)=9t-\frac{15}{\pi}\sin (\frac{\pi t}{5})$. Since $\sin0=\sin(2\pi)=0$,
\[\overline h=\frac{1}{10}\int_0^{10}h(t)\,\dd t=\frac{F(10)-F(0)}{10}=\frac{90}{10}=9\ \mathrm{m}.\]

\item \textit{[0.5 pt]} The second derivative is $h''(t)=\frac{3\pi^2}{25}\cos (\frac{\pi t}{5})$. Its sign is that of the cosine: positive on $[0,\frac{5}{2}[$ and $]\frac{15}{2},10]$, negative between these two intervals. The function is convex on $[0,\frac{5}{2}]$ and $[\frac{15}{2},10]$, concave on $[\frac{5}{2},\frac{15}{2}]$. The two changes of convexity give the points of inflection $(\frac{5}{2},9)$ and $(\frac{15}{2},9)$.

\item \textit{[0.5 pt]} Between $\frac{10}{3}$ and $\frac{20}{3}$; duration $\frac{10}{3}$ s.

\item \textit{[0.5 pt]} $h'$ reaches $\frac{3}{5}\pi$ $\mathrm{m\,s^{-1}}$ at $t=\frac{5}{2}$ and its opposite at $t=\frac{15}{2}$.

\item \textit{[1 pt]} Displacement $h(10)-h(0)=0$. Rise of $6$ m then fall of $6$ m: total distance $12$ m, equal to the integral of $|h'|$.

\end{enumerate}

\subsection*{Exercise 3 · Line, plane and projection}\begin{enumerate}

\item \textit{[0.5 pt]} $(x,y,z)=(2+t,2t,2t)$, $t\in\R$.

\item \textit{[0.5 pt]} $2+9t=11$ gives $t=1$; $H=(3,2,2)$.

\item \textit{[0.5 pt]} The vector $\vec n=(1,2,2)$ is normal to $\mathcal P$. Since $H\in\mathcal P$ and $\vect{MH}=\vec n$, $H$ is the orthogonal projection of $M$ onto the plane. The coordinate system being orthonormal, $\|\vec n\|=\sqrt{1+4+4}=3$. Hence
\[d(M,\mathcal P)=MH=|1|\,3=3.\]

\item \textit{[1 pt]} Set $y=z=s$; $x=11-4s$. Hence $(x,y,z)=(11,0,0)+s(-4,1,1)$.

\item \textit{[0.5 pt]} Yes, its last two coordinates are equal. $H$ lies on the line and $\vect{MH}$ is orthogonal to every direction of the plane, hence to this line. Distance $3$.

\item \textit{[0.5 pt]} $x+2y+2z=2$; same normal vector and passing through $M$.

\item \textit{[0.5 pt]} $\vec n\cdot(-4,1,1)=-4+2+2=0$: angle $\frac{\pi}{2}$.

\item \textit{[1 pt]} $N=M+\vec n$ lies on the same normal line as $M$; its projection onto $\mathcal P$ is therefore also $H$. Now $H\in\mathcal P\cap\mathcal Q$. For every $P$ on this line, $NP^2=NH^2+HP^2$, so its minimal distance to $N$ is $NH$. Finally,
\[d(N,\mathcal P)=d(N,\mathcal P\cap\mathcal Q)=\|\vect{NH}\|=|1-1|\,\|\vec n\|=0.\]

\end{enumerate}

\subsection*{Exercise 4 · Repetitions and estimation}\begin{enumerate}

\item \textit{[0.5 pt]} Repetition of a fixed number $20$ of identical independent trials with two outcomes; $X\sim\mathcal B(20,\frac{1}{5})$.

\item \textit{[0.5 pt]} $P(X=2)=\binom{20}2(\frac{1}{5})^2(\frac{4}{5})^{18}\simeq0.136909$.

\item \textit{[0.5 pt]} $1-P(X=0)-P(X=1)=1-(\frac{4}{5})^{20}-20(\frac{1}{5})(\frac{4}{5})^{19}\simeq0.930825$.

\item \textit{[1 pt]} $E(X)=4$, $V(X)=\frac{16}{5}$.

\item \textit{[0.5 pt]} For $N$ independent trials, the probability of no success is $(1-p)^N$. The required probability is therefore $1-(\frac{4}{5})^N$. It is strictly increasing in $N$, since $0<\frac{4}{5}<1$. The calculations give
\[P_{13}\simeq0.945024<0.95,\qquad P_{14}\simeq0.956020\ge0.95.\]
Thus $14$ is the smallest suitable number of trials.

\item \textit{[0.5 pt]} Let $n=20$. For $j\in\{0,\ldots,n\}$, an ordered outcome containing $j$ successes and $n-j$ failures has, by independence, probability $p^j(1-p)^{n-j}$. There are $\binom nj$ choices for the positions of the successes. These outcomes are mutually exclusive, so
\[P(X=j)=\binom nj p^j(1-p)^{n-j},\qquad p=\frac{1}{5}.\]

\item \textit{[0.5 pt]} If $B_1,\ldots,B_N$ are the independent indicator variables of the successes, $F_N=\frac{B_1+\cdots+B_N}{N}$. With $E(B_i)=p$ and $V(B_i)=p(1-p)$,
\[E(F_N)=\frac{Np}{N}=\frac{1}{5},\qquad V(F_N)=\frac{Np(1-p)}{N^2}=\frac{\frac{4}{25}}{N}.\]
Independence is required for adding variances, not for the linearity of expected value.

\item \textit{[1 pt]} The Bienaymé–Chebyshev inequality applied to $F_N$ gives
\[P(|F_N-p|\ge0.02)\le\frac{p(1-p)}{N(0.02)^2}.\]
To guarantee that this probability is at most $0.05$, it suffices that
\[N\ge\frac{\frac{4}{25}}{0.05\times(0.02)^2}=8000.\]
Taking the complement, we then have $P(|F_N-p|<0.02)\ge0.95$. The size $8000$ is sufficient; the inequality being an upper bound, it does not necessarily give the minimum size for the exact probability.

\end{enumerate}

\fi

\clearpage\phantomsection\label{sujet:9-2}\addcontentsline{toc}{chapter}{Final mock · Practice 2}

\section*{Final mock · Practice 2}

{\sffamily\bfseries\color{teal}PRACTICE PAPER · TRAINING\par}

\textbf{Duration: 4 h · Marked out of 20 · Four independent exercises.}\\Calculator allowed. Every answer must be justified; incomplete working will be taken into account. Give exact values; round probabilities to $10^{-6}$ and other approximations to $10^{-4}$, unless stated otherwise.

\smallskip\textbf{Name:}\enspace\dotfill\quad\textbf{Class:}\enspace\dotfill\quad\textbf{Date:}\enspace\dotfill

\subsection*{Exercise 1 · A logarithmic function \hfill 5 points}Consider the function $f$ defined on $]0;+\infty[$ by $f(x)=\frac{5+\ln x}x$. Let $\mathcal C_f$ denote its graph.

\begin{enumerate}

\item Determine the limits of $f$ at $0^+$ and at $+\infty$. \hfill\textit{(0.5 pt)}

\item Calculate $f' (x)$, then draw up the table of variations of $f$. \hfill\textit{(0.5 pt)}

\item State the maximum value of $f$, then solve $f(x)=0$. \hfill\textit{(0.5 pt)}

\item Study the convexity of $f$ and determine the coordinates of the point of inflection of $\mathcal C_f$. \hfill\textit{(1 pt)}

\item Prove that the equation $f(x)=\frac12 e^{4}$ has exactly two solutions in $]0;+\infty[$. \hfill\textit{(0.5 pt)}

\item Find an antiderivative of $f$ on $]0;+\infty[$. \hfill\textit{(0.5 pt)}

\item Calculate the mean value of $f$ on $[1;e]$. \hfill\textit{(0.5 pt)}

\item Independently of the previous questions, solve $\ln(x-3)+\ln(x+3)=\ln(40)$, after stating the set on which this equation is defined. \hfill\textit{(1 pt)}

\end{enumerate}

\subsection*{Exercise 2 · Periodic phenomenon \hfill 5 points}A point moves on a wheel turning at constant speed in a vertical plane. Its height above the ground is modelled by $h(t)=12-4\cos\left(\frac{\pi t}{4}\right)$, where $t$ is expressed in seconds and $h(t)$ in metres. We study one complete revolution, for $t\in[0;8]$.

\begin{enumerate}

\item Give a period of the model extended to $\R$, then the minimum and maximum heights over one revolution. \hfill\textit{(0.5 pt)}

\item Calculate $h'(t)$ and draw up the table of variations of $h$ on $[0;8]$. \hfill\textit{(0.5 pt)}

\item Find all times $t\in[0;8]$ such that $h(t)=14$. \hfill\textit{(0.5 pt)}

\item Calculate the mean value of $h$ on $[0;8]$. \hfill\textit{(1 pt)}

\item Study the convexity of $h$ on the interval of study and determine the coordinates of the points of inflection of its graph. \hfill\textit{(0.5 pt)}

\item Determine the total time during which the height is at least equal to the level in question 3. \hfill\textit{(0.5 pt)}

\item The vertical speed is given by $h'(t)$. Determine its maximum and minimum values, in metres per second, and the times at which they are reached. \hfill\textit{(0.5 pt)}

\item Calculate the vertical displacement $h(8)-h(0)$, then the sum of the height travelled on the way up and that travelled on the way down. Explain why these two results are different. \hfill\textit{(1 pt)}

\end{enumerate}

\clearpage\noindent{\sffamily\bfseries\color{navy}Final mock · Practice 2 · continued}\par

\subsection*{Exercise 3 · Line, plane and projection \hfill 5 points}In an orthonormal coordinate system of space, consider the plane $\mathcal P$ with equation $x+2y+2z=7$ and the point $M(3;0;0)$. The line $d$ passes through $M$ and has direction vector $\vec n=(1;2;2)$.

\begin{enumerate}

\item Give a parametric representation of $d$. \hfill\textit{(0.5 pt)}

\item Prove that $d$ meets $\mathcal P$ at a point $H$, then calculate the coordinates of $H$. \hfill\textit{(0.5 pt)}

\item Justify that $H$ is the orthogonal projection of $M$ onto $\mathcal P$, then calculate the distance from $M$ to $\mathcal P$. \hfill\textit{(0.5 pt)}

\item Consider the plane $\mathcal Q$ with equation $y=z$. Prove that $\mathcal P\cap\mathcal Q$ is a line and give a parametric representation of it. \hfill\textit{(1 pt)}

\item Check that $H\in\mathcal Q$. Deduce the distance from $M$ to the line $\mathcal P\cap\mathcal Q$. \hfill\textit{(0.5 pt)}

\item Determine a Cartesian equation of the plane parallel to $\mathcal P$ passing through $M$. \hfill\textit{(0.5 pt)}

\item Determine the angle between $d$ and the line $\mathcal P\cap\mathcal Q$. \hfill\textit{(0.5 pt)}

\item Define the point $N$ by $\vect{MN}=\vec n$. Calculate the distances from $N$ to the plane $\mathcal P$ and to the line $\mathcal P\cap\mathcal Q$, then compare them. \hfill\textit{(1 pt)}

\end{enumerate}

\subsection*{Exercise 4 · Repetitions and estimation \hfill 5 points}We repeat $17$ times a trial with two outcomes, called success and failure. The trials are independent and the probability of success is $p=\frac{3}{10}$ on each trial. Let $X$ denote the number of successes obtained.

\begin{enumerate}

\item Justify that $X$ follows a binomial distribution and state its parameters. \hfill\textit{(0.5 pt)}

\item Calculate the probability of obtaining exactly two successes. \hfill\textit{(0.5 pt)}

\item Using complementary events, express and then calculate the probability of obtaining at least two successes. \hfill\textit{(0.5 pt)}

\item Calculate the expected value and the variance of $X$. \hfill\textit{(1 pt)}

\item We now repeat the same trial a number $N\ge1$ of times, still independently. Determine the smallest integer $N$ for which the probability of obtaining at least one success is greater than or equal to $0.95$. A search using a table or a loop is allowed. \hfill\textit{(0.5 pt)}

\item For an integer $j\in\{0,\ldots,17\}$, prove the formula for $P(X=j)$ by counting the possible positions of the successes. \hfill\textit{(0.5 pt)}

\item For $N$ independent trials with the same parameter $p$, let $F_N$ denote the frequency of successes. Prove the expressions for $E(F_N)$ and $V(F_N)$ in terms of $N$ and $p$. \hfill\textit{(0.5 pt)}

\item Using the Bienaymé–Chebyshev inequality, determine a sufficient size $N$ so that $P(|F_N-p|<0.02)\ge0.95$. Is the size obtained necessarily the smallest possible for this probability? Justify. \hfill\textit{(1 pt)}

\end{enumerate}

\ifprof\clearpage

\section*{Solutions · Final mock · Practice 2}\label{cor:9-2}

\subsection*{Exercise 1 · A logarithmic function}\begin{enumerate}

\item \textit{[0.5 pt]} At $0^+$, the numerator tends to $-\infty$ and the denominator to $0^+$: limit $-\infty$; at $+\infty$, limit $0$ by comparison of growth rates.

\item \textit{[0.5 pt]} By the quotient rule,
\[f'(x)=\frac{(\frac{1}{x})x-(5+\ln x)}{x^2}=\frac{-4-\ln x}{x^2}.\]
The denominator is positive. The numerator vanishes at $s=e^{-4}$, is positive for $0<x<s$ and negative for $x>s$. Hence $f$ is strictly increasing on $]0,s]$, then strictly decreasing on $[s,+\infty[$.

\item \textit{[0.5 pt]} The maximum is attained only at $s=e^{-4}$ and equals
\[f(s)=\frac{5- 4}{e^{-4}}=e^{4}.\]
Since $x>0$, $f(x)=0$ is equivalent to $5+\ln x=0$, that is $x=e^{-5}$.

\item \textit{[1 pt]} Differentiating $(-4-\ln x)x^{-2}$:
\[f''(x)=-x^{-3}-2(-4-\ln x)x^{-3}=\frac{2\ln x+7}{x^3}.\]
The numerator vanishes at $t=e^{-\frac{7}{2}}$ and changes from negative to positive. The function is concave on $]0,t]$, convex on $[t,+\infty[$. It changes convexity at $t$ and $f(t)=\frac{3}{2t}$: the point of inflection is $I(t,\frac{3}{2t})$.

\item \textit{[0.5 pt]} Let $m=\frac{e^{4}}{2}$ and $s=e^{-4}$. The function is continuous on $]0,+\infty[$. On $]0,s]$, it is strictly increasing, with limit $-\infty$ at $0$ and value $2m$ at $s$: the corollary of the IVT gives exactly one solution in $]0,s[$. On $[s,+\infty[$, it is strictly decreasing from $2m$ to $0<m$: exactly one solution in $]s,+\infty[$. These intervals cover the domain and $f(s)\ne m$: there are exactly two solutions.

\item \textit{[0.5 pt]} Let $F(x)=5\ln x+\frac12(\ln x)^2$, differentiable for $x>0$. Then $F'(x)=\frac{5}{x}+\frac{\ln x}{x}=f(x)$; $F$ is an antiderivative on $]0,+\infty[$.

\item \textit{[0.5 pt]} With this antiderivative, $F(e)=\frac{11}{2}$ and $F(1)=0$. The integral therefore equals $\frac{11}{2}$, and the mean value is
\[\mu=\frac{1}{e-1}\int_1^e f(x)\,\dd x=\frac{11}{2(e-1)}.\]

\item \textit{[1 pt]} Domain $x>3$; $x^2-9=40$, so $x=7$, the negative root being excluded.

\end{enumerate}

\subsection*{Exercise 2 · Periodic phenomenon}\begin{enumerate}

\item \textit{[0.5 pt]} Period $8$ s; minimum $8$ m, maximum $16$ m.

\item \textit{[0.5 pt]} $h'(t)=\pi\sin (\frac{\pi t}{4})$; increasing on $[0,4]$, decreasing on $[4,8]$.

\item \textit{[0.5 pt]} The equation is equivalent to $\cos (\frac{\pi t}{4})=-\frac{1}{2}$. Over one revolution, $\frac{\pi t}{4}\in[0,2\pi]$; the suitable angles are $\frac{2\pi}{3}$ and $\frac{4\pi}{3}$. Hence $t=\frac{8}{3}$ s or $t=\frac{16}{3}$ s.

\item \textit{[1 pt]} An antiderivative is $F(t)=12t-\frac{16}{\pi}\sin (\frac{\pi t}{4})$. Since $\sin0=\sin(2\pi)=0$,
\[\overline h=\frac{1}{8}\int_0^{8}h(t)\,\dd t=\frac{F(8)-F(0)}{8}=\frac{96}{8}=12\ \mathrm{m}.\]

\item \textit{[0.5 pt]} The second derivative is $h''(t)=\frac{4\pi^2}{16}\cos (\frac{\pi t}{4})$. Its sign is that of the cosine: positive on $[0,2[$ and $]6,8]$, negative between these two intervals. The function is convex on $[0,2]$ and $[6,8]$, concave on $[2,6]$. The two changes of convexity give the points of inflection $(2,12)$ and $(6,12)$.

\item \textit{[0.5 pt]} Between $\frac{8}{3}$ and $\frac{16}{3}$; duration $\frac{8}{3}$ s.

\item \textit{[0.5 pt]} $h'$ reaches $1\pi$ $\mathrm{m\,s^{-1}}$ at $t=2$ and its opposite at $t=6$.

\item \textit{[1 pt]} Displacement $h(8)-h(0)=0$. Rise of $8$ m then fall of $8$ m: total distance $16$ m, equal to the integral of $|h'|$.

\end{enumerate}

\subsection*{Exercise 3 · Line, plane and projection}\begin{enumerate}

\item \textit{[0.5 pt]} $(x,y,z)=(3+t,2t,2t)$, $t\in\R$.

\item \textit{[0.5 pt]} $3+9t=7$ gives $t=\frac{4}{9}$; $H=(\frac{31}{9},\frac{8}{9},\frac{8}{9})$.

\item \textit{[0.5 pt]} The vector $\vec n=(1,2,2)$ is normal to $\mathcal P$. Since $H\in\mathcal P$ and $\vect{MH}=\frac{4}{9}\vec n$, $H$ is the orthogonal projection of $M$ onto the plane. The coordinate system being orthonormal, $\|\vec n\|=\sqrt{1+4+4}=3$. Hence
\[d(M,\mathcal P)=MH=|\frac{4}{9}|\,3=\frac{4}{3}.\]

\item \textit{[1 pt]} Set $y=z=s$; $x=7-4s$. Hence $(x,y,z)=(7,0,0)+s(-4,1,1)$.

\item \textit{[0.5 pt]} Yes, its last two coordinates are equal. $H$ lies on the line and $\vect{MH}$ is orthogonal to every direction of the plane, hence to this line. Distance $\frac{4}{3}$.

\item \textit{[0.5 pt]} $x+2y+2z=3$; same normal vector and passing through $M$.

\item \textit{[0.5 pt]} $\vec n\cdot(-4,1,1)=-4+2+2=0$: angle $\frac{\pi}{2}$.

\item \textit{[1 pt]} $N=M+\vec n$ lies on the same normal line as $M$; its projection onto $\mathcal P$ is therefore also $H$. Now $H\in\mathcal P\cap\mathcal Q$. For every $P$ on this line, $NP^2=NH^2+HP^2$, so its minimal distance to $N$ is $NH$. Finally,
\[d(N,\mathcal P)=d(N,\mathcal P\cap\mathcal Q)=\|\vect{NH}\|=|\frac{4}{9}-1|\,\|\vec n\|=\frac{5}{3}.\]

\end{enumerate}

\subsection*{Exercise 4 · Repetitions and estimation}\begin{enumerate}

\item \textit{[0.5 pt]} Repetition of a fixed number $17$ of identical independent trials with two outcomes; $X\sim\mathcal B(17,\frac{3}{10})$.

\item \textit{[0.5 pt]} $P(X=2)=\binom{17}2(\frac{3}{10})^2(\frac{7}{10})^{15}\simeq0.058110$.

\item \textit{[0.5 pt]} $1-P(X=0)-P(X=1)=1-(\frac{7}{10})^{17}-17(\frac{3}{10})(\frac{7}{10})^{16}\simeq0.980725$.

\item \textit{[1 pt]} $E(X)=\frac{51}{10}$, $V(X)=\frac{357}{100}$.

\item \textit{[0.5 pt]} For $N$ independent trials, the probability of no success is $(1-p)^N$. The required probability is therefore $1-(\frac{7}{10})^N$. It is strictly increasing in $N$, since $0<\frac{7}{10}<1$. The calculations give
\[P_{8}\simeq0.942352<0.95,\qquad P_{9}\simeq0.959646\ge0.95.\]
Thus $9$ is the smallest suitable number of trials.

\item \textit{[0.5 pt]} Let $n=17$. For $j\in\{0,\ldots,n\}$, an ordered outcome containing $j$ successes and $n-j$ failures has, by independence, probability $p^j(1-p)^{n-j}$. There are $\binom nj$ choices for the positions of the successes. These outcomes are mutually exclusive, so
\[P(X=j)=\binom nj p^j(1-p)^{n-j},\qquad p=\frac{3}{10}.\]

\item \textit{[0.5 pt]} If $B_1,\ldots,B_N$ are the independent indicator variables of the successes, $F_N=\frac{B_1+\cdots+B_N}{N}$. With $E(B_i)=p$ and $V(B_i)=p(1-p)$,
\[E(F_N)=\frac{Np}{N}=\frac{3}{10},\qquad V(F_N)=\frac{Np(1-p)}{N^2}=\frac{\frac{21}{100}}{N}.\]
Independence is required for adding variances, not for the linearity of expected value.

\item \textit{[1 pt]} The Bienaymé–Chebyshev inequality applied to $F_N$ gives
\[P(|F_N-p|\ge0.02)\le\frac{p(1-p)}{N(0.02)^2}.\]
To guarantee that this probability is at most $0.05$, it suffices that
\[N\ge\frac{\frac{21}{100}}{0.05\times(0.02)^2}=10500.\]
Taking the complement, we then have $P(|F_N-p|<0.02)\ge0.95$. The size $10500$ is sufficient; the inequality being an upper bound, it does not necessarily give the minimum size for the exact probability.

\end{enumerate}

\fi

\ifprof

\clearpage\phantomsection\label{sujet:9-3}\addcontentsline{toc}{chapter}{Final mock · A}

\section*{Final mock · A}

{\sffamily\bfseries\color{amber}PAPER RESERVED FOR ASSESSMENT\par}

\textbf{Duration: 4 h · Marked out of 20 · Four independent exercises.}\\Calculator allowed. Every answer must be justified; incomplete working will be taken into account. Give exact values; round probabilities to $10^{-6}$ and other approximations to $10^{-4}$, unless stated otherwise.

\smallskip\textbf{Name:}\enspace\dotfill\quad\textbf{Class:}\enspace\dotfill\quad\textbf{Date:}\enspace\dotfill

\subsection*{Exercise 1 · A logarithmic function \hfill 5 points}Consider the function $f$ defined on $]0;+\infty[$ by $f(x)=\frac{4+\ln x}x$. Let $\mathcal C_f$ denote its graph.

\begin{enumerate}

\item Determine the limits of $f$ at $0^+$ and at $+\infty$. \hfill\textit{(0.5 pt)}

\item Calculate $f' (x)$, then draw up the table of variations of $f$. \hfill\textit{(0.5 pt)}

\item State the maximum value of $f$, then solve $f(x)=0$. \hfill\textit{(0.5 pt)}

\item Study the convexity of $f$ and determine the coordinates of the point of inflection of $\mathcal C_f$. \hfill\textit{(1 pt)}

\item Prove that the equation $f(x)=\frac12 e^{3}$ has exactly two solutions in $]0;+\infty[$. \hfill\textit{(0.5 pt)}

\item Find an antiderivative of $f$ on $]0;+\infty[$. \hfill\textit{(0.5 pt)}

\item Calculate the mean value of $f$ on $[1;e]$. \hfill\textit{(0.5 pt)}

\item Independently of the previous questions, solve $\ln(x-1)+\ln(x+1)=\ln(15)$, after stating the set on which this equation is defined. \hfill\textit{(1 pt)}

\end{enumerate}

\subsection*{Exercise 2 · Periodic phenomenon \hfill 5 points}A point moves on a wheel turning at constant speed in a vertical plane. Its height above the ground is modelled by $h(t)=7-2\cos\left(\frac{\pi t}{3}\right)$, where $t$ is expressed in seconds and $h(t)$ in metres. We study one complete revolution, for $t\in[0;6]$.

\begin{enumerate}

\item Give a period of the model extended to $\R$, then the minimum and maximum heights over one revolution. \hfill\textit{(0.5 pt)}

\item Calculate $h'(t)$ and draw up the table of variations of $h$ on $[0;6]$. \hfill\textit{(0.5 pt)}

\item Find all times $t\in[0;6]$ such that $h(t)=8$. \hfill\textit{(0.5 pt)}

\item Calculate the mean value of $h$ on $[0;6]$. \hfill\textit{(1 pt)}

\item Study the convexity of $h$ on the interval of study and determine the coordinates of the points of inflection of its graph. \hfill\textit{(0.5 pt)}

\item Determine the total time during which the height is at least equal to the level in question 3. \hfill\textit{(0.5 pt)}

\item The vertical speed is given by $h'(t)$. Determine its maximum and minimum values, in metres per second, and the times at which they are reached. \hfill\textit{(0.5 pt)}

\item Calculate the vertical displacement $h(6)-h(0)$, then the sum of the height travelled on the way up and that travelled on the way down. Explain why these two results are different. \hfill\textit{(1 pt)}

\end{enumerate}

\clearpage\noindent{\sffamily\bfseries\color{navy}Final mock · A · continued}\par

\subsection*{Exercise 3 · Line, plane and projection \hfill 5 points}In an orthonormal coordinate system of space, consider the plane $\mathcal P$ with equation $x+2y+2z=8$ and the point $M(1;0;0)$. The line $d$ passes through $M$ and has direction vector $\vec n=(1;2;2)$.

\begin{enumerate}

\item Give a parametric representation of $d$. \hfill\textit{(0.5 pt)}

\item Prove that $d$ meets $\mathcal P$ at a point $H$, then calculate the coordinates of $H$. \hfill\textit{(0.5 pt)}

\item Justify that $H$ is the orthogonal projection of $M$ onto $\mathcal P$, then calculate the distance from $M$ to $\mathcal P$. \hfill\textit{(0.5 pt)}

\item Consider the plane $\mathcal Q$ with equation $y=z$. Prove that $\mathcal P\cap\mathcal Q$ is a line and give a parametric representation of it. \hfill\textit{(1 pt)}

\item Check that $H\in\mathcal Q$. Deduce the distance from $M$ to the line $\mathcal P\cap\mathcal Q$. \hfill\textit{(0.5 pt)}

\item Determine a Cartesian equation of the plane parallel to $\mathcal P$ passing through $M$. \hfill\textit{(0.5 pt)}

\item Determine the angle between $d$ and the line $\mathcal P\cap\mathcal Q$. \hfill\textit{(0.5 pt)}

\item Define the point $N$ by $\vect{MN}=\vec n$. Calculate the distances from $N$ to the plane $\mathcal P$ and to the line $\mathcal P\cap\mathcal Q$, then compare them. \hfill\textit{(1 pt)}

\end{enumerate}

\subsection*{Exercise 4 · Repetitions and estimation \hfill 5 points}We repeat $18$ times a trial with two outcomes, called success and failure. The trials are independent and the probability of success is $p=\frac{1}{10}$ on each trial. Let $X$ denote the number of successes obtained.

\begin{enumerate}

\item Justify that $X$ follows a binomial distribution and state its parameters. \hfill\textit{(0.5 pt)}

\item Calculate the probability of obtaining exactly two successes. \hfill\textit{(0.5 pt)}

\item Using complementary events, express and then calculate the probability of obtaining at least two successes. \hfill\textit{(0.5 pt)}

\item Calculate the expected value and the variance of $X$. \hfill\textit{(1 pt)}

\item We now repeat the same trial a number $N\ge1$ of times, still independently. Determine the smallest integer $N$ for which the probability of obtaining at least one success is greater than or equal to $0.95$. A search using a table or a loop is allowed. \hfill\textit{(0.5 pt)}

\item For an integer $j\in\{0,\ldots,18\}$, prove the formula for $P(X=j)$ by counting the possible positions of the successes. \hfill\textit{(0.5 pt)}

\item For $N$ independent trials with the same parameter $p$, let $F_N$ denote the frequency of successes. Prove the expressions for $E(F_N)$ and $V(F_N)$ in terms of $N$ and $p$. \hfill\textit{(0.5 pt)}

\item Using the Bienaymé–Chebyshev inequality, determine a sufficient size $N$ so that $P(|F_N-p|<0.02)\ge0.95$. Is the size obtained necessarily the smallest possible for this probability? Justify. \hfill\textit{(1 pt)}

\end{enumerate}

\ifprof\clearpage

\section*{Solutions · Final mock · A}\label{cor:9-3}

\subsection*{Exercise 1 · A logarithmic function}\begin{enumerate}

\item \textit{[0.5 pt]} At $0^+$, the numerator tends to $-\infty$ and the denominator to $0^+$: limit $-\infty$; at $+\infty$, limit $0$ by comparison of growth rates.

\item \textit{[0.5 pt]} By the quotient rule,
\[f'(x)=\frac{(\frac{1}{x})x-(4+\ln x)}{x^2}=\frac{-3-\ln x}{x^2}.\]
The denominator is positive. The numerator vanishes at $s=e^{-3}$, is positive for $0<x<s$ and negative for $x>s$. Hence $f$ is strictly increasing on $]0,s]$, then strictly decreasing on $[s,+\infty[$.

\item \textit{[0.5 pt]} The maximum is attained only at $s=e^{-3}$ and equals
\[f(s)=\frac{4- 3}{e^{-3}}=e^{3}.\]
Since $x>0$, $f(x)=0$ is equivalent to $4+\ln x=0$, that is $x=e^{-4}$.

\item \textit{[1 pt]} Differentiating $(-3-\ln x)x^{-2}$:
\[f''(x)=-x^{-3}-2(-3-\ln x)x^{-3}=\frac{2\ln x+5}{x^3}.\]
The numerator vanishes at $t=e^{-\frac{5}{2}}$ and changes from negative to positive. The function is concave on $]0,t]$, convex on $[t,+\infty[$. It changes convexity at $t$ and $f(t)=\frac{3}{2t}$: the point of inflection is $I(t,\frac{3}{2t})$.

\item \textit{[0.5 pt]} Let $m=\frac{e^{3}}{2}$ and $s=e^{-3}$. The function is continuous on $]0,+\infty[$. On $]0,s]$, it is strictly increasing, with limit $-\infty$ at $0$ and value $2m$ at $s$: the corollary of the IVT gives exactly one solution in $]0,s[$. On $[s,+\infty[$, it is strictly decreasing from $2m$ to $0<m$: exactly one solution in $]s,+\infty[$. These intervals cover the domain and $f(s)\ne m$: there are exactly two solutions.

\item \textit{[0.5 pt]} Let $F(x)=4\ln x+\frac12(\ln x)^2$, differentiable for $x>0$. Then $F'(x)=\frac{4}{x}+\frac{\ln x}{x}=f(x)$; $F$ is an antiderivative on $]0,+\infty[$.

\item \textit{[0.5 pt]} With this antiderivative, $F(e)=\frac{9}{2}$ and $F(1)=0$. The integral therefore equals $\frac{9}{2}$, and the mean value is
\[\mu=\frac{1}{e-1}\int_1^e f(x)\,\dd x=\frac{9}{2(e-1)}.\]

\item \textit{[1 pt]} Domain $x>1$; $x^2-1=15$, so $x=4$, the negative root being excluded.

\end{enumerate}

\subsection*{Exercise 2 · Periodic phenomenon}\begin{enumerate}

\item \textit{[0.5 pt]} Period $6$ s; minimum $5$ m, maximum $9$ m.

\item \textit{[0.5 pt]} $h'(t)=(\frac{2}{3}\pi)\sin (\frac{\pi t}{3})$; increasing on $[0,3]$, decreasing on $[3,6]$.

\item \textit{[0.5 pt]} The equation is equivalent to $\cos (\frac{\pi t}{3})=-\frac{1}{2}$. Over one revolution, $\frac{\pi t}{3}\in[0,2\pi]$; the suitable angles are $\frac{2\pi}{3}$ and $\frac{4\pi}{3}$. Hence $t=2$ s or $t=4$ s.

\item \textit{[1 pt]} An antiderivative is $F(t)=7t-\frac{6}{\pi}\sin (\frac{\pi t}{3})$. Since $\sin0=\sin(2\pi)=0$,
\[\overline h=\frac{1}{6}\int_0^{6}h(t)\,\dd t=\frac{F(6)-F(0)}{6}=\frac{42}{6}=7\ \mathrm{m}.\]

\item \textit{[0.5 pt]} The second derivative is $h''(t)=\frac{2\pi^2}{9}\cos (\frac{\pi t}{3})$. Its sign is that of the cosine: positive on $[0,\frac{3}{2}[$ and $]\frac{9}{2},6]$, negative between these two intervals. The function is convex on $[0,\frac{3}{2}]$ and $[\frac{9}{2},6]$, concave on $[\frac{3}{2},\frac{9}{2}]$. The two changes of convexity give the points of inflection $(\frac{3}{2},7)$ and $(\frac{9}{2},7)$.

\item \textit{[0.5 pt]} Between $2$ and $4$; duration $2$ s.

\item \textit{[0.5 pt]} $h'$ reaches $\frac{2}{3}\pi$ $\mathrm{m\,s^{-1}}$ at $t=\frac{3}{2}$ and its opposite at $t=\frac{9}{2}$.

\item \textit{[1 pt]} Displacement $h(6)-h(0)=0$. Rise of $4$ m then fall of $4$ m: total distance $8$ m, equal to the integral of $|h'|$.

\end{enumerate}

\subsection*{Exercise 3 · Line, plane and projection}\begin{enumerate}

\item \textit{[0.5 pt]} $(x,y,z)=(1+t,2t,2t)$, $t\in\R$.

\item \textit{[0.5 pt]} $1+9t=8$ gives $t=\frac{7}{9}$; $H=(\frac{16}{9},\frac{14}{9},\frac{14}{9})$.

\item \textit{[0.5 pt]} The vector $\vec n=(1,2,2)$ is normal to $\mathcal P$. Since $H\in\mathcal P$ and $\vect{MH}=\frac{7}{9}\vec n$, $H$ is the orthogonal projection of $M$ onto the plane. The coordinate system being orthonormal, $\|\vec n\|=\sqrt{1+4+4}=3$. Hence
\[d(M,\mathcal P)=MH=|\frac{7}{9}|\,3=\frac{7}{3}.\]

\item \textit{[1 pt]} Set $y=z=s$; $x=8-4s$. Hence $(x,y,z)=(8,0,0)+s(-4,1,1)$.

\item \textit{[0.5 pt]} Yes, its last two coordinates are equal. $H$ lies on the line and $\vect{MH}$ is orthogonal to every direction of the plane, hence to this line. Distance $\frac{7}{3}$.

\item \textit{[0.5 pt]} $x+2y+2z=1$; same normal vector and passing through $M$.

\item \textit{[0.5 pt]} $\vec n\cdot(-4,1,1)=-4+2+2=0$: angle $\frac{\pi}{2}$.

\item \textit{[1 pt]} $N=M+\vec n$ lies on the same normal line as $M$; its projection onto $\mathcal P$ is therefore also $H$. Now $H\in\mathcal P\cap\mathcal Q$. For every $P$ on this line, $NP^2=NH^2+HP^2$, so its minimal distance to $N$ is $NH$. Finally,
\[d(N,\mathcal P)=d(N,\mathcal P\cap\mathcal Q)=\|\vect{NH}\|=|\frac{7}{9}-1|\,\|\vec n\|=\frac{2}{3}.\]

\end{enumerate}

\subsection*{Exercise 4 · Repetitions and estimation}\begin{enumerate}

\item \textit{[0.5 pt]} Repetition of a fixed number $18$ of identical independent trials with two outcomes; $X\sim\mathcal B(18,\frac{1}{10})$.

\item \textit{[0.5 pt]} $P(X=2)=\binom{18}2(\frac{1}{10})^2(\frac{9}{10})^{16}\simeq0.283512$.

\item \textit{[0.5 pt]} $1-P(X=0)-P(X=1)=1-(\frac{9}{10})^{18}-18(\frac{1}{10})(\frac{9}{10})^{17}\simeq0.549716$.

\item \textit{[1 pt]} $E(X)=\frac{9}{5}$, $V(X)=\frac{81}{50}$.

\item \textit{[0.5 pt]} For $N$ independent trials, the probability of no success is $(1-p)^N$. The required probability is therefore $1-(\frac{9}{10})^N$. It is strictly increasing in $N$, since $0<\frac{9}{10}<1$. The calculations give
\[P_{28}\simeq0.947665<0.95,\qquad P_{29}\simeq0.952899\ge0.95.\]
Thus $29$ is the smallest suitable number of trials.

\item \textit{[0.5 pt]} Let $n=18$. For $j\in\{0,\ldots,n\}$, an ordered outcome containing $j$ successes and $n-j$ failures has, by independence, probability $p^j(1-p)^{n-j}$. There are $\binom nj$ choices for the positions of the successes. These outcomes are mutually exclusive, so
\[P(X=j)=\binom nj p^j(1-p)^{n-j},\qquad p=\frac{1}{10}.\]

\item \textit{[0.5 pt]} If $B_1,\ldots,B_N$ are the independent indicator variables of the successes, $F_N=\frac{B_1+\cdots+B_N}{N}$. With $E(B_i)=p$ and $V(B_i)=p(1-p)$,
\[E(F_N)=\frac{Np}{N}=\frac{1}{10},\qquad V(F_N)=\frac{Np(1-p)}{N^2}=\frac{\frac{9}{100}}{N}.\]
Independence is required for adding variances, not for the linearity of expected value.

\item \textit{[1 pt]} The Bienaymé–Chebyshev inequality applied to $F_N$ gives
\[P(|F_N-p|\ge0.02)\le\frac{p(1-p)}{N(0.02)^2}.\]
To guarantee that this probability is at most $0.05$, it suffices that
\[N\ge\frac{\frac{9}{100}}{0.05\times(0.02)^2}=4500.\]
Taking the complement, we then have $P(|F_N-p|<0.02)\ge0.95$. The size $4500$ is sufficient; the inequality being an upper bound, it does not necessarily give the minimum size for the exact probability.

\end{enumerate}

\fi

\fi

\ifprof

\clearpage\phantomsection\label{sujet:9-4}\addcontentsline{toc}{chapter}{Final mock · B}

\section*{Final mock · B}

{\sffamily\bfseries\color{amber}PAPER RESERVED FOR ASSESSMENT\par}

\textbf{Duration: 4 h · Marked out of 20 · Four independent exercises.}\\Calculator allowed. Every answer must be justified; incomplete working will be taken into account. Give exact values; round probabilities to $10^{-6}$ and other approximations to $10^{-4}$, unless stated otherwise.

\smallskip\textbf{Name:}\enspace\dotfill\quad\textbf{Class:}\enspace\dotfill\quad\textbf{Date:}\enspace\dotfill

\subsection*{Exercise 1 · A logarithmic function \hfill 5 points}Consider the function $f$ defined on $]0;+\infty[$ by $f(x)=\frac{3+\ln x}x$. Let $\mathcal C_f$ denote its graph.

\begin{enumerate}

\item Determine the limits of $f$ at $0^+$ and at $+\infty$. \hfill\textit{(0.5 pt)}

\item Calculate $f' (x)$, then draw up the table of variations of $f$. \hfill\textit{(0.5 pt)}

\item State the maximum value of $f$, then solve $f(x)=0$. \hfill\textit{(0.5 pt)}

\item Study the convexity of $f$ and determine the coordinates of the point of inflection of $\mathcal C_f$. \hfill\textit{(1 pt)}

\item Prove that the equation $f(x)=\frac12 e^{2}$ has exactly two solutions in $]0;+\infty[$. \hfill\textit{(0.5 pt)}

\item Find an antiderivative of $f$ on $]0;+\infty[$. \hfill\textit{(0.5 pt)}

\item Calculate the mean value of $f$ on $[1;e]$. \hfill\textit{(0.5 pt)}

\item Independently of the previous questions, solve $\ln(x-2)+\ln(x+2)=\ln(32)$, after stating the set on which this equation is defined. \hfill\textit{(1 pt)}

\end{enumerate}

\subsection*{Exercise 2 · Periodic phenomenon \hfill 5 points}A point moves on a wheel turning at constant speed in a vertical plane. Its height above the ground is modelled by $h(t)=10-3\cos\left(\frac{\pi t}{2}\right)$, where $t$ is expressed in seconds and $h(t)$ in metres. We study one complete revolution, for $t\in[0;4]$.

\begin{enumerate}

\item Give a period of the model extended to $\R$, then the minimum and maximum heights over one revolution. \hfill\textit{(0.5 pt)}

\item Calculate $h'(t)$ and draw up the table of variations of $h$ on $[0;4]$. \hfill\textit{(0.5 pt)}

\item Find all times $t\in[0;4]$ such that $h(t)=\frac{23}{2}$. \hfill\textit{(0.5 pt)}

\item Calculate the mean value of $h$ on $[0;4]$. \hfill\textit{(1 pt)}

\item Study the convexity of $h$ on the interval of study and determine the coordinates of the points of inflection of its graph. \hfill\textit{(0.5 pt)}

\item Determine the total time during which the height is at least equal to the level in question 3. \hfill\textit{(0.5 pt)}

\item The vertical speed is given by $h'(t)$. Determine its maximum and minimum values, in metres per second, and the times at which they are reached. \hfill\textit{(0.5 pt)}

\item Calculate the vertical displacement $h(4)-h(0)$, then the sum of the height travelled on the way up and that travelled on the way down. Explain why these two results are different. \hfill\textit{(1 pt)}

\end{enumerate}

\clearpage\noindent{\sffamily\bfseries\color{navy}Final mock · B · continued}\par

\subsection*{Exercise 3 · Line, plane and projection \hfill 5 points}In an orthonormal coordinate system of space, consider the plane $\mathcal P$ with equation $x+2y+2z=12$ and the point $M(2;0;0)$. The line $d$ passes through $M$ and has direction vector $\vec n=(1;2;2)$.

\begin{enumerate}

\item Give a parametric representation of $d$. \hfill\textit{(0.5 pt)}

\item Prove that $d$ meets $\mathcal P$ at a point $H$, then calculate the coordinates of $H$. \hfill\textit{(0.5 pt)}

\item Justify that $H$ is the orthogonal projection of $M$ onto $\mathcal P$, then calculate the distance from $M$ to $\mathcal P$. \hfill\textit{(0.5 pt)}

\item Consider the plane $\mathcal Q$ with equation $y=z$. Prove that $\mathcal P\cap\mathcal Q$ is a line and give a parametric representation of it. \hfill\textit{(1 pt)}

\item Check that $H\in\mathcal Q$. Deduce the distance from $M$ to the line $\mathcal P\cap\mathcal Q$. \hfill\textit{(0.5 pt)}

\item Determine a Cartesian equation of the plane parallel to $\mathcal P$ passing through $M$. \hfill\textit{(0.5 pt)}

\item Determine the angle between $d$ and the line $\mathcal P\cap\mathcal Q$. \hfill\textit{(0.5 pt)}

\item Define the point $N$ by $\vect{MN}=\vec n$. Calculate the distances from $N$ to the plane $\mathcal P$ and to the line $\mathcal P\cap\mathcal Q$, then compare them. \hfill\textit{(1 pt)}

\end{enumerate}

\subsection*{Exercise 4 · Repetitions and estimation \hfill 5 points}We repeat $20$ times a trial with two outcomes, called success and failure. The trials are independent and the probability of success is $p=\frac{3}{10}$ on each trial. Let $X$ denote the number of successes obtained.

\begin{enumerate}

\item Justify that $X$ follows a binomial distribution and state its parameters. \hfill\textit{(0.5 pt)}

\item Calculate the probability of obtaining exactly two successes. \hfill\textit{(0.5 pt)}

\item Using complementary events, express and then calculate the probability of obtaining at least two successes. \hfill\textit{(0.5 pt)}

\item Calculate the expected value and the variance of $X$. \hfill\textit{(1 pt)}

\item We now repeat the same trial a number $N\ge1$ of times, still independently. Determine the smallest integer $N$ for which the probability of obtaining at least one success is greater than or equal to $0.95$. A search using a table or a loop is allowed. \hfill\textit{(0.5 pt)}

\item For an integer $j\in\{0,\ldots,20\}$, prove the formula for $P(X=j)$ by counting the possible positions of the successes. \hfill\textit{(0.5 pt)}

\item For $N$ independent trials with the same parameter $p$, let $F_N$ denote the frequency of successes. Prove the expressions for $E(F_N)$ and $V(F_N)$ in terms of $N$ and $p$. \hfill\textit{(0.5 pt)}

\item Using the Bienaymé–Chebyshev inequality, determine a sufficient size $N$ so that $P(|F_N-p|<0.02)\ge0.95$. Is the size obtained necessarily the smallest possible for this probability? Justify. \hfill\textit{(1 pt)}

\end{enumerate}

\ifprof\clearpage

\section*{Solutions · Final mock · B}\label{cor:9-4}

\subsection*{Exercise 1 · A logarithmic function}\begin{enumerate}

\item \textit{[0.5 pt]} At $0^+$, the numerator tends to $-\infty$ and the denominator to $0^+$: limit $-\infty$; at $+\infty$, limit $0$ by comparison of growth rates.

\item \textit{[0.5 pt]} By the quotient rule,
\[f'(x)=\frac{(\frac{1}{x})x-(3+\ln x)}{x^2}=\frac{-2-\ln x}{x^2}.\]
The denominator is positive. The numerator vanishes at $s=e^{-2}$, is positive for $0<x<s$ and negative for $x>s$. Hence $f$ is strictly increasing on $]0,s]$, then strictly decreasing on $[s,+\infty[$.

\item \textit{[0.5 pt]} The maximum is attained only at $s=e^{-2}$ and equals
\[f(s)=\frac{3- 2}{e^{-2}}=e^{2}.\]
Since $x>0$, $f(x)=0$ is equivalent to $3+\ln x=0$, that is $x=e^{-3}$.

\item \textit{[1 pt]} Differentiating $(-2-\ln x)x^{-2}$:
\[f''(x)=-x^{-3}-2(-2-\ln x)x^{-3}=\frac{2\ln x+3}{x^3}.\]
The numerator vanishes at $t=e^{-\frac{3}{2}}$ and changes from negative to positive. The function is concave on $]0,t]$, convex on $[t,+\infty[$. It changes convexity at $t$ and $f(t)=\frac{3}{2t}$: the point of inflection is $I(t,\frac{3}{2t})$.

\item \textit{[0.5 pt]} Let $m=\frac{e^{2}}{2}$ and $s=e^{-2}$. The function is continuous on $]0,+\infty[$. On $]0,s]$, it is strictly increasing, with limit $-\infty$ at $0$ and value $2m$ at $s$: the corollary of the IVT gives exactly one solution in $]0,s[$. On $[s,+\infty[$, it is strictly decreasing from $2m$ to $0<m$: exactly one solution in $]s,+\infty[$. These intervals cover the domain and $f(s)\ne m$: there are exactly two solutions.

\item \textit{[0.5 pt]} Let $F(x)=3\ln x+\frac12(\ln x)^2$, differentiable for $x>0$. Then $F'(x)=\frac{3}{x}+\frac{\ln x}{x}=f(x)$; $F$ is an antiderivative on $]0,+\infty[$.

\item \textit{[0.5 pt]} With this antiderivative, $F(e)=\frac{7}{2}$ and $F(1)=0$. The integral therefore equals $\frac{7}{2}$, and the mean value is
\[\mu=\frac{1}{e-1}\int_1^e f(x)\,\dd x=\frac{7}{2(e-1)}.\]

\item \textit{[1 pt]} Domain $x>2$; $x^2-4=32$, so $x=6$, the negative root being excluded.

\end{enumerate}

\subsection*{Exercise 2 · Periodic phenomenon}\begin{enumerate}

\item \textit{[0.5 pt]} Period $4$ s; minimum $7$ m, maximum $13$ m.

\item \textit{[0.5 pt]} $h'(t)=(\frac{3}{2}\pi)\sin (\frac{\pi t}{2})$; increasing on $[0,2]$, decreasing on $[2,4]$.

\item \textit{[0.5 pt]} The equation is equivalent to $\cos (\frac{\pi t}{2})=-\frac{1}{2}$. Over one revolution, $\frac{\pi t}{2}\in[0,2\pi]$; the suitable angles are $\frac{2\pi}{3}$ and $\frac{4\pi}{3}$. Hence $t=\frac{4}{3}$ s or $t=\frac{8}{3}$ s.

\item \textit{[1 pt]} An antiderivative is $F(t)=10t-\frac{6}{\pi}\sin (\frac{\pi t}{2})$. Since $\sin0=\sin(2\pi)=0$,
\[\overline h=\frac{1}{4}\int_0^{4}h(t)\,\dd t=\frac{F(4)-F(0)}{4}=\frac{40}{4}=10\ \mathrm{m}.\]

\item \textit{[0.5 pt]} The second derivative is $h''(t)=\frac{3\pi^2}{4}\cos (\frac{\pi t}{2})$. Its sign is that of the cosine: positive on $[0,1[$ and $]3,4]$, negative between these two intervals. The function is convex on $[0,1]$ and $[3,4]$, concave on $[1,3]$. The two changes of convexity give the points of inflection $(1,10)$ and $(3,10)$.

\item \textit{[0.5 pt]} Between $\frac{4}{3}$ and $\frac{8}{3}$; duration $\frac{4}{3}$ s.

\item \textit{[0.5 pt]} $h'$ reaches $\frac{3}{2}\pi$ $\mathrm{m\,s^{-1}}$ at $t=1$ and its opposite at $t=3$.

\item \textit{[1 pt]} Displacement $h(4)-h(0)=0$. Rise of $6$ m then fall of $6$ m: total distance $12$ m, equal to the integral of $|h'|$.

\end{enumerate}

\subsection*{Exercise 3 · Line, plane and projection}\begin{enumerate}

\item \textit{[0.5 pt]} $(x,y,z)=(2+t,2t,2t)$, $t\in\R$.

\item \textit{[0.5 pt]} $2+9t=12$ gives $t=\frac{10}{9}$; $H=(\frac{28}{9},\frac{20}{9},\frac{20}{9})$.

\item \textit{[0.5 pt]} The vector $\vec n=(1,2,2)$ is normal to $\mathcal P$. Since $H\in\mathcal P$ and $\vect{MH}=\frac{10}{9}\vec n$, $H$ is the orthogonal projection of $M$ onto the plane. The coordinate system being orthonormal, $\|\vec n\|=\sqrt{1+4+4}=3$. Hence
\[d(M,\mathcal P)=MH=|\frac{10}{9}|\,3=\frac{10}{3}.\]

\item \textit{[1 pt]} Set $y=z=s$; $x=12-4s$. Hence $(x,y,z)=(12,0,0)+s(-4,1,1)$.

\item \textit{[0.5 pt]} Yes, its last two coordinates are equal. $H$ lies on the line and $\vect{MH}$ is orthogonal to every direction of the plane, hence to this line. Distance $\frac{10}{3}$.

\item \textit{[0.5 pt]} $x+2y+2z=2$; same normal vector and passing through $M$.

\item \textit{[0.5 pt]} $\vec n\cdot(-4,1,1)=-4+2+2=0$: angle $\frac{\pi}{2}$.

\item \textit{[1 pt]} $N=M+\vec n$ lies on the same normal line as $M$; its projection onto $\mathcal P$ is therefore also $H$. Now $H\in\mathcal P\cap\mathcal Q$. For every $P$ on this line, $NP^2=NH^2+HP^2$, so its minimal distance to $N$ is $NH$. Finally,
\[d(N,\mathcal P)=d(N,\mathcal P\cap\mathcal Q)=\|\vect{NH}\|=|\frac{10}{9}-1|\,\|\vec n\|=\frac{1}{3}.\]

\end{enumerate}

\subsection*{Exercise 4 · Repetitions and estimation}\begin{enumerate}

\item \textit{[0.5 pt]} Repetition of a fixed number $20$ of identical independent trials with two outcomes; $X\sim\mathcal B(20,\frac{3}{10})$.

\item \textit{[0.5 pt]} $P(X=2)=\binom{20}2(\frac{3}{10})^2(\frac{7}{10})^{18}\simeq0.027846$.

\item \textit{[0.5 pt]} $1-P(X=0)-P(X=1)=1-(\frac{7}{10})^{20}-20(\frac{3}{10})(\frac{7}{10})^{19}\simeq0.992363$.

\item \textit{[1 pt]} $E(X)=6$, $V(X)=\frac{21}{5}$.

\item \textit{[0.5 pt]} For $N$ independent trials, the probability of no success is $(1-p)^N$. The required probability is therefore $1-(\frac{7}{10})^N$. It is strictly increasing in $N$, since $0<\frac{7}{10}<1$. The calculations give
\[P_{8}\simeq0.942352<0.95,\qquad P_{9}\simeq0.959646\ge0.95.\]
Thus $9$ is the smallest suitable number of trials.

\item \textit{[0.5 pt]} Let $n=20$. For $j\in\{0,\ldots,n\}$, an ordered outcome containing $j$ successes and $n-j$ failures has, by independence, probability $p^j(1-p)^{n-j}$. There are $\binom nj$ choices for the positions of the successes. These outcomes are mutually exclusive, so
\[P(X=j)=\binom nj p^j(1-p)^{n-j},\qquad p=\frac{3}{10}.\]

\item \textit{[0.5 pt]} If $B_1,\ldots,B_N$ are the independent indicator variables of the successes, $F_N=\frac{B_1+\cdots+B_N}{N}$. With $E(B_i)=p$ and $V(B_i)=p(1-p)$,
\[E(F_N)=\frac{Np}{N}=\frac{3}{10},\qquad V(F_N)=\frac{Np(1-p)}{N^2}=\frac{\frac{21}{100}}{N}.\]
Independence is required for adding variances, not for the linearity of expected value.

\item \textit{[1 pt]} The Bienaymé–Chebyshev inequality applied to $F_N$ gives
\[P(|F_N-p|\ge0.02)\le\frac{p(1-p)}{N(0.02)^2}.\]
To guarantee that this probability is at most $0.05$, it suffices that
\[N\ge\frac{\frac{21}{100}}{0.05\times(0.02)^2}=10500.\]
Taking the complement, we then have $P(|F_N-p|<0.02)\ge0.95$. The size $10500$ is sufficient; the inequality being an upper bound, it does not necessarily give the minimum size for the exact probability.

\end{enumerate}

\fi

\fi

